\documentclass[10pt]{article}
\usepackage[margin=20mm]{geometry}
\usepackage{amsmath,amssymb,fontspec,unicode-math,xurl,hyperref}
\IfFontExistsTF{DejaVu Serif}{\setmainfont{DejaVu Serif}}{\setmainfont{Latin Modern Roman}}
\setmathfont{latinmodern-math.otf}
\setlength{\parindent}{0pt}\setlength{\parskip}{6pt}\sloppy
\begin{document}
\section*{Scientific verification report: I25-40163 (0056)}
Inspection outcome: No Issue Found within the inspected scope; two scope clarifications retained. Author adoption has not been requested.\par
\subsection*{1. Identity and reviewed version}
Patrik Reizinger, Randall Balestriero, David Klindt and Wieland Brendel, Position: An Empirically Grounded Identifiability Theory Will Accelerate Self-Supervised Learning Research. ICML 2025, PMLR 267, 82072-82090 [E1-E2]. The reviewed object is the 19-page publisher-linked PDF, not an assumed arXiv equivalent. Locations below use one-based physical PDF pages, which match the printed article pages.\par
SHA-256 (verified against the acquisition receipt, with the original bytes preserved):\par
`f749eddc389ae66806ec82a77742bf0575dd880fa1ff4d906510615966130ee2`\par
\subsection*{2. Scope and decisive result}
The review read the argument in Sections 1-6 (pp. 1-9), Table 1, Figure 1, and Appendices A-C (pp. 18-19). Selected bibliography entries were traced; this is not an audit of every cited paper. The central checks concern the move from absolute to relative identifiability, the conditional interpretation of the Platonic Representation Hypothesis (PRH), and whether a proposed research programme is represented as a completed theorem.\par
The decisive counterevidence to an overly strong reading is in the paper itself: Section 2 says representation similarity is not universal; Section 4 explicitly presents Singular Identifiability Theory (SITh) as a future blueprint; and Appendix B distinguishes similarity between learned models from recovery of a common ground truth. Different augmentations can select different latent factors. The paper therefore does not claim a theorem that every SSL architecture, dataset or training run converges to one universal representation [E1, pp. 2, 4, 7, 18-19].\par
No demonstrated central mathematical or empirical error was established in these checks. This conclusion does not certify all cited results or every broad historical assertion.\par
\subsection*{3. Checked reasoning and operative assumptions}
Conditional identifiability bridge. Suppose two learned representations refer to the same latent vector z, the same population, and the same dimension, and satisfy f1(x) = A1 z and f2(x) = A2 z with invertible square matrices A1 and A2. Substitution gives f2(x) = A2 A1\textasciicircum{}-1 f1(x). This is an analytic identity, not a finite-training result. With a fixed equivalence relation, the same implication follows from symmetry and transitivity. It fails as an argument for arbitrary models if the latent variables, admissible transformations or populations differ. Appendix B recognizes these restrictions rather than discharging them universally [E1, p. 18].\par
Upstream SimCLR result. Section 3.2’s example was compared with Zimmermann et al.’s published Sections 3.1-3.2, Theorems 1-2 [E3]. Their spherical result uses a uniform latent marginal, a von Mises-Fisher conditional, a differentiable injective generative map, a matching conditional model, population loss minimization and an asymptotic number of negative samples. It recovers latents up to an orthogonal linear transformation. The position paper uses this as an idealized example and explicitly identifies the practical gap. This supports its conditional illustration; it does not establish finite-data recovery for deployed SimCLR.\par
Evidence from Figure 1. The rendered p. 7 was inspected. Its caption specifies RBF-kernel CKA with sigma = 0.4 and attributes the similarity matrix to Ciernik et al.; the right panel is a UMAP depiction. The text acknowledges non-universal similarity. Neither the matrix nor a UMAP projection is a proof of identifiability. The report does not substitute a linear-CKA counterexample for a claim about this different metric. The underlying embeddings and plotting pipeline were not reproduced.\par
Research agenda and evaluation. Questions 4.1.1-4.8.2, including finite samples, optimization, architecture, projection and compositional generalization, are presented as questions. A theorem proving that SITh will accelerate research is not supplied or assumed. The critique of using one classification score to assess all latent information is logically coherent: a task insensitive to a factor cannot by itself establish recovery of that factor [E1, pp. 4-8].\par
\subsection*{4. Clarifications and unresolved boundaries}
C1: Fix the admissible relation in Definition A.2. Literally, the existence of an unrestricted equivalence relation is vacuous: the relation that makes every pair equivalent always satisfies the definition. A nonvacuous reading fixes a scientifically meaningful transformation class before quantifying over learned models. The surrounding discussion repeatedly uses simple or linear transformations, while A.1 refers to a particular relation. No downstream proof was found to depend on the unrestricted reading. This is retained as a definition/scope clarification, without a supported defect severity. A precise clarification would begin: “For a pre-specified admissible equivalence relation …”. This is an editorial interpretation, not a verified repair of a false theorem [E1, p. 18].\par
C2: Identifiability is not universal metric agreement. Invertible linear equivalence alone does not guarantee high values for every similarity measure or recovery after arbitrary finite training. The paper’s conditional wording and open questions are material counterevidence against treating such a stronger statement as its claim. No counterexample to the actual statement is asserted.\par
Pending source-level boundary. Appendix B’s statement about non-content factors collapsing needs to be read with the representation dimension, layer and augmentation assumptions of von Kuegelgen et al. Section 4.3.1 itself discusses other factors being learned with larger representations and more training. This review did not reconstruct that upstream theorem in full; it records the boundary as Needs Evidence, without inferring a contradiction or increasing severity [E1, pp. 5, 18].\par
\subsection*{5. Dependence, versions and remaining work}
The checked dependency is: idealized identifiability results can explain representation agreement under shared assumptions; their practical gaps motivate SITh. The conditional bridge survives. The empirical effectiveness of the proposed future programme remains an open research outcome, not a reproduced result.\par
The arXiv record identifies 2504.13101v3, revised 24 July 2025, as an ICML camera-ready version [E4]. Its PDF was not substituted for the publisher object or compared byte for byte. No separate erratum was identified on the inspected publisher/arXiv pages; a complete correction-history search was not performed. OpenReview returned a browser-verification page, so its discussion and attachments were not inspected. No paper-specific executable implementation was linked on the inspected publisher page, and the paper does not introduce a new implementation requiring execution. Author code was not run.\par
Fresh experiments would be needed to reproduce Figure 1, compare optimizers or architectures, or test whether the proposed agenda improves downstream outcomes. Those experiments are outside this static review.\par
\subsection*{6. Conclusion table}
Item | Issue judgment | Repair verification | Adoption / remaining work\par
Conditional PRH / identifiability reasoning | No Issue Found in inspected scope | Not applicable | Human decision pending\par
A.2 admissible relation | Scope clarification | Restricted wording justified analytically; no manuscript repair applied | Authors not contacted\par
Collapse wording / upstream assumptions | Needs Evidence, no supported severity | Not performed | Reconstruct cited theorem if pursued\par
Figure 1 empirical provenance | Inspected caption and argument; not reproduced | Not applicable | Original embeddings/code needed for reproduction\par
\subsection*{7. Consulted sources}
- E1. Reizinger et al. (2025), exact publisher PDF, `papers/0056/0056.pdf`; extraction: `sources/0056/paper.layout.txt`.\par
- E2. Publisher record: https://proceedings.mlr.press/v267/reizinger25a.html ; snapshot: `sources/0056/publisher.html`.\par
- E3. Zimmermann, Sharma, Schneider, Bethge and Brendel (2021), Contrastive Learning Inverts the Data Generating Process, PMLR 139:12979-12990. https://proceedings.mlr.press/v139/zimmermann21a.html ; published PDF and extraction stored in `sources/0056/`. Only the theorem statements and their immediate setup were checked, not its complete proof.\par
- E4. https://arxiv.org/abs/2504.13101 ; version-history snapshot: `sources/0056/arxiv\_record.html`.\par
- Workflow reference. Kassis, Agarwal, He, Patel and Brueckner (2026), Scientific Agent Skills: A Library of Procedural Knowledge for Research Agents. https://doi.org/10.48550/arXiv.2609.00065 . Metadata checked against arXiv; v2 listed on 2 September 2026. This describes review procedure, not evidence for the paper’s scientific claims.\par
\subsection*{Supporting checks}
Identity and SHA-256 were checked before scientific analysis.\par
No human verification, reviewer endorsement or author approval is claimed. The original paper and source-acquisition record were not modified.\par
\end{document}
